\section{El punto final: ¿dejarán de importar los datos?}

Al final de la entrevista se pregunta: ¿llegará el día en que el problema de los datos deje de importar por completo? Xie Chen (谢晨) responde que cada vez está más convencido de que, cuanto más fuerte es la inteligencia, más hambre tiene de conocimiento y de datos. Cuanto más capaz es una persona, más quiere aprender, y con la IA podría ocurrir lo mismo. El desenlace quizá no sea que «los datos desaparezcan», sino pasar de aprender de datos externos a aprender de forma continua en simulación y en entornos propios.

\subsection{La evaluación es el cuello de botella}

En el caso de los robots, Xie Chen considera que el problema más importante podría ser escalar la evaluación. Sin una evaluación a escala no se sabe si el modelo de verdad se volvió más inteligente. Los modelos de lenguaje grandes y los Agent también necesitan métricas de evaluación de nivel superior, porque cuanto más fuerte es el modelo, más fácil es que los benchmark anteriores dejen de servir.

\begin{importantbox}{La evaluación también es un problema de datos}
El conjunto de evaluación no es un juez externo, sino una fuente de retroalimentación para el entrenamiento y la investigación. Sin una buena evaluación, el ciclo cerrado de datos no puede saber hacia dónde avanzar.
\end{importantbox}

\subsection{Autoaprendizaje y el desenlace en simulación}

Si la IA puede, dentro de un entorno de simulación, fijarse objetivos, intentar, fallar, corregir y volver a intentar, la fábrica de datos no desaparecerá, sino que cambiará de forma: pasará de recolectar datos humanos a construir entornos, objetivos y retroalimentación para que la IA aprenda en ellos de forma continua.

\begin{knowledgebox}{De la data factory a la environment factory}
Cuando el rendimiento marginal de los datos externos disminuye, lo que de verdad importa podría ser construir entornos de los que se pueda aprender. La simulación, la evaluación, la reward y la retroalimentación de los fallos se convertirán en la nueva fábrica de datos.
\end{knowledgebox}

\subsection{Resumen de la sección}

Los datos no desaparecerán: pasarán de ser un corpus estático a ser entornos dinámicos. Cuanto más fuerte es la inteligencia, más necesita retroalimentación de mayor calidad y entornos de aprendizaje más complejos.
